\bookchapter{บทที่ 8: การสื่อสาร: ข้อความ โทรศัพท์ และโซเชียล}

\chapterintro{การสื่อสารคือภารกิจดั้งเดิมและเป็นแกนกลางที่สุดของโทรศัพท์ ทุกวันนี้มันวิวัฒนาการจากการโทรและส่งข้อความธรรมดา กลายเป็นระบบนิเวศที่ซับซ้อนซึ่งรวมการสื่อสารแบบทันที เครือข่ายสังคม และอีเมลเข้าไว้ด้วยกัน}

\booksection{ฟังก์ชันขั้นสูงของแอปโทรศัพท์}

แอปโทรศัพท์ในมือถือสมัยนี้ไปไกลกว่าแป้นหมุนเลขธรรมดาแล้ว มันรวมเครื่องมือทรงพลังอย่างการบันทึกเสียง การบล็อกสายรบกวน และการระบุข้อมูลอัจฉริยะไว้ในที่เดียว

\booksubsection{การบันทึกเสียงสนทนา (โปรดใส่ใจเรื่องความชอบด้วยกฎหมาย):}

\begin{itemize}
  \item ประโยชน์: บันทึกเนื้อหาการสนทนาที่สำคัญ ใช้สำหรับบันทึกการประชุม สัมภาษณ์ คำสัญญาจากลูกค้า หรือเก็บไว้เป็นหลักฐานเมื่อมีข้อพิพาท
  \item ข้อกฎหมาย: ในจีน การใช้เสียงบันทึกสนทนาเป็นพยานหลักฐานมีเงื่อนไขว่าต้องได้รับความยินยอมจากอีกฝ่าย หรือใช้เพื่อปกป้องสิทธิ์และผลประโยชน์ที่ชอบด้วยกฎหมายของตนเอง
  \item การบันทึกฝ่ายเดียว เทียบกับความยินยอมสองฝ่าย: ในทางกฎหมาย เสียงบันทึกที่ได้มา "โดยไม่ได้รับความยินยอมจากอีกฝ่าย" ไม่ได้เป็นโมฆะโดยอัตโนมัติ ตามคำตีความตุลาการที่เกี่ยวข้องของศาลประชาชนสูงสุด ตราบใดที่วิธีที่ได้เสียงบันทึกมานั้นไม่ละเมิดสิทธิ์และผลประโยชน์ของผู้อื่นอย่างร้ายแรง ไม่ขัดต่อกฎหมายห้าม และไม่ขัดต่อความสงบเรียบร้อยและศีลธรรมอันดี ก็ยังสามารถยื่นเป็นพยานหลักฐานได้ อย่างไรก็ตาม ศาลในบางมณฑลเข้มงวดกว่าในการรับฟัง "การแอบบันทึก" ในคดีแพ่ง ส่วนคดีอาญามักต้องตรวจสอบความชอบด้วยกฎหมายเข้มงวดกว่า
  \item ความแตกต่างตามพื้นที่: ศาลในปักกิ่ง เซี่ยงไฮ้ กวางตุ้ง และที่อื่น ๆ ค่อนข้างผ่อนปรนต่อการใช้เสียงบันทึกฝ่ายเดียวเป็นพยานหลักฐานทางแพ่ง แต่เสียงบันทึกที่เกี่ยวกับความลับทางการค้าหรือความเป็นส่วนตัว ไม่ว่าจะอยู่มณฑลใดก็ต้องระวังเป็นพิเศษ เพราะอาจถูกตัดออกเนื่องจากละเมิดสิทธิ์ความเป็นส่วนตัว
  \item วิธีใช้งาน: บน Android มักมีปุ่ม "บันทึกเสียง" ชัดเจนบนหน้าจอสนทนา ส่วนระบบ iOS ไม่มีฟังก์ชันบันทึกเสียงสนทนาในตัว ต้องใช้แอปบุคคลที่สาม หรือใช้หูฟัง Bluetooth ที่มีฟังก์ชันบันทึกเสียง
\end{itemize}

\begin{tcbwarn}{เส้นแดงในการใช้เสียงบันทึก}
\begin{itemize}
  \item ห้ามใช้เสียงบันทึกเพื่อแบล็กเมล์ เผยแพร่ หรือดูหมิ่นผู้อื่นในที่สาธารณะ
  \item เนื้อหาที่เกี่ยวข้องกับความเป็นส่วนตัวของบุคคลที่สาม (เช่น บังเอิญอัดเสียงสนทนาของคนอื่น) ห้ามเผยแพร่โดยพลการ
  \item ในสถานการณ์สำคัญ แนะนำให้แจ้งด้วยวาจาตั้งแต่ต้นการสนทนาว่า "การสนทนานี้จะถูกบันทึกเสียง" ทั้งถูกต้องตามกฎและมีน้ำหนักเป็นพยานหลักฐานมากขึ้น
\end{itemize}

\end{tcbwarn}

\booksubsection{เนื้อหาเชิงลึกเรื่องการจัดการสายรบกวน:}

\begin{itemize}
  \item หลักการ: แอปโทรศัพท์ค้นฐานข้อมูลเบอร์ขนาดใหญ่ผ่านเครือข่าย เพื่อแสดง "สายรบกวน" "สายหลอกลวง" หรือชื่อบริษัทขณะมีสายเข้า
  \item วิธีเปิดใช้งาน: บน Android ไปที่แอปโทรศัพท์ > การตั้งค่า > การบล็อกสายรบกวน บน iOS ไปที่ การตั้งค่า > โทรศัพท์ > การบล็อกและการระบุตัวตนสายเข้า สำหรับเส้นทางเปิดปิดและตัวเลือกที่แนะนำของแต่ละแบรนด์ ดูตารางเปรียบเทียบแบรนด์ในบทที่ 13 "เครื่องมือระบบใช้งานจริง"
  \item แพลตฟอร์มติดป้ายหลัก: การติดป้ายเบอร์ในจีนให้บริการโดยผู้ให้บริการข้อมูลสามรายหลัก ได้แก่ Tencent Mobile Manager, 360 Mobile Guard และ Tapir (Tapir) ผู้ให้บริการเครือข่าย (China Mobile/China Unicom/China Telecom) ก็มักเรียกใช้ข้อมูลเหล่านี้สำหรับการแสดง "ชื่อเจ้าของเครื่อง" และฟังก์ชัน "บล็อกสายรบกวน" เช่นกัน ฐานข้อมูลติดป้ายของผู้ผลิตมือถือ (Huawei, Xiaomi, OPPO, vivo) ส่วนใหญ่ก็ร่วมมือกับผู้ให้บริการข้างต้น
  \item ถ้าถูกติดป้ายผิดพลาดทำอย่างไร: หากเบอร์ทำงานปกติของคุณถูกติดป้ายผิดว่าเป็น "อาหารเดลิเวอรี" "พัสดุ" หรือ "สายรบกวน" สามารถอุทธรณ์ผ่านช่องทางต่อไปนี้: Tapir: เข้าไปที่แพลตฟอร์มอุทธรณ์เบอร์ (haoma.tk.cn) เพื่อยื่นเอกสารยืนยันตัวตน
  \item 360: ยื่นผ่านฟังก์ชัน "อุทธรณ์เบอร์" ในแอป 360 Mobile Guard
  \item Tencent: ยื่นผ่านหน้า "อุทธรณ์เบอร์" ในมินิโปรแกรม "Tencent Mobile Manager"
  \item ผู้ให้บริการเครือข่ายรายใหญ่สามราย: ผู้ใช้ China Mobile สามารถส่งข้อความ "KTFSR" ไปที่ 10086 เพื่อเปิดฟังก์ชันป้องกันสายรบกวน ผู้ใช้ China Telecom ติดตามบัญชีทางการ "Tianyi Anti-Harassment" ผู้ใช้ China Unicom ตั้งค่าใน "Unicom Mobile Manager"
\end{itemize}

ขั้นสูงเรื่องป้องกันสายรบกวน: เปิด "การบล็อกสายเรียกเข้าความถี่สูง" ของผู้ให้บริการเครือข่าย (เช่น การจำกัดสายเรียกเข้าความถี่สูงของ China Mobile) จะสามารถบล็อกสายโทรออกด้วยเสียง AI โดยอัตโนมัติ

ความสามารถบล็อกทั่วไปที่มือถือมีในตัว: การระบุสายเข้า / การบล็อกอัจฉริยะ: ติดป้ายและบล็อกสายรบกวน สายหลอกลวง สายโฆษณาขายของโดยอัตโนมัติตามฐานข้อมูลเบอร์

บัญชีดำ / บัญชีขาว: บล็อกเบอร์รบกวนที่รู้จักด้วยตนเอง เพิ่มรายชื่อติดต่อ เบอร์พัสดุและอาหารเดลิเวอรีเข้าบัญชีขาว เพื่อไม่ให้บล็อกผิด

ตัวกรองคำสำคัญในข้อความ: ข้อความที่มีคำสำคัญเช่น "ถูกรางวัล" "กู้ยืม" "คลิกลิงก์" จะถูกจัดเข้ากล่องบล็อกโดยอัตโนมัติ

iOS ปิดเสียงสายจากคนแปลกหน้า: สายเข้าที่ไม่ได้อยู่ในรายชื่อติดต่อจะถูกปิดเสียงเข้าวอยซ์เมลเท่านั้น ไม่รบกวนคุณ เหมาะกับผู้ที่ได้รับสายรบกวนบ่อย แต่ต้องระวังสายสำคัญจากคนแปลกหน้าที่อาจพลาด

เชื่อมโยงกับโหมดห้ามรบกวน: ตอนกลางคืนหรือระหว่างประชุม อนุญาตเฉพาะสายจากรายชื่อติดต่อเท่านั้น จับคู่กับการตั้งเวลาเปิดปิด คุณภาพการนอนจะดีขึ้นทันที

ใช้แอปช่วยอย่างไรให้เข้ากัน: ใช้ตัวบล็อกของระบบเป็นหลักก่อน แล้วเสริมด้วย National Anti-Fraud Center และการป้องกันสายรบกวนของผู้ให้บริการเครือข่าย สำหรับแอปทำความสะอาดสามารถเลือกเครื่องมือเบา ๆ อย่าง Qingzhuo แต่ติดตั้งแอปความปลอดภัย/ทำความสะอาดประเภทเดียวกันเพียงชุดเดียวก็พอ ดูรายละเอียดในบทที่ 13

\booksubsection{VoLTE/VoNR คืออะไร และเปิดใช้อย่างไร:}

\begin{itemize}
  \item VoLTE (Voice over LTE): เทคโนโลยีเสียงสนทนาบนเครือข่าย 4G ขณะโทรจะไม่ตกลงไปใช้ 2G/3G อีก สามารถเล่นเน็ตขณะโทรได้ คุณภาพเสียงชัดเจนกว่าการโทรแบบเดิม และติดต่อได้เร็วกว่า
  \item VoNR (Voice over NR): เสียงสนทนาบน 5G เป็นวิวัฒนาการรุ่นถัดไปของ VoLTE หน่วงเวลาน้อยกว่า คุณภาพเสียงดีกว่า และรองรับวิดีโอคอลความละเอียดสูง ต้องมีครบทั้งมือถือ 5G + แพ็กเกจ 5G + พื้นที่ที่มีสัญญาณ 5G จึงจะใช้ได้
  \item วิธีเปิดใช้งาน: บน iOS เปิด "VoLTE" ที่ "การตั้งค่า > เครือข่ายมือถือ > ตัวเลือกข้อมูลมือถือ > เสียงและข้อมูล" บน Android หาสวิตช์ "VoLTE คุณภาพเสียงสูง" ที่ "การตั้งค่า > เครือข่ายมือถือ > ข้อมูลมือถือ > ขั้นสูง" ผู้ให้บริการเครือข่ายมักเปิดใช้เป็นค่าเริ่มต้น หากพบว่าอินเทอร์เน็ตตัดขณะโทร หรือเล่นเน็ตไม่ได้ ส่วนใหญ่เป็นเพราะยังไม่ได้เปิด VoLTE
\end{itemize}

\begin{tcbinfo}{เกร็ดน่ารู้}
หากมือถือของคุณอินเทอร์เน็ตช้าลงหรือตัดกะทันหันขณะโทร ส่วนใหญ่เป็นเพราะยังไม่ได้เปิด VoLTE/VoNR จึงตกลงไปใช้เครือข่าย 2G/3G เมื่อเปิดแล้วจะสามารถ “โทรไปเล่นวิดีโอไป” ได้

\end{tcbinfo}

\booksection{ข้อความและ iMessage/RCS}

ข้อความ SMS แบบเดิมยังไม่ล้าสมัย ยังคงเป็นช่องทางที่เชื่อถือได้ในการรับรหัสยืนยัน การแจ้งเตือนจากธนาคาร และข้อมูลทางการ

\begin{itemize}
  \item iMessage: ข้อความที่ส่งผ่านอินเทอร์เน็ตระหว่างอุปกรณ์ Apple แสดงเป็นฟองสีน้ำเงิน รองรับรูปภาพวิดีโอคุณภาพสูง แจ้งเตือนอ่านแล้ว และอื่น ๆ ฟรีทั้งหมดและเข้ารหัสต้นทางถึงปลายทาง
  \item RCS: โปรโตคอลการสื่อสารรุ่นถัดไปที่ผลักดันโดยผู้ให้บริการเครือข่ายทั่วโลกและ Google ฟังก์ชันคล้าย iMessage หรือ WeChat แต่อิงตามหมายเลขโทรศัพท์
  \item การทำงานร่วมกันของ iMessage และ RCS: ตั้งแต่ iOS 18 Apple เริ่มรองรับมาตรฐาน RCS ซึ่งหมายความว่าสามารถส่งรูปภาพความละเอียดสูง แจ้งเตือนอ่านแล้ว และสถานะกำลังพิมพ์ระหว่าง iPhone กับมือถือ Android ผ่าน RCS ได้ หมดยุคประสบการณ์ "ฟองสีเขียว" ที่ขาดตอน แต่เงื่อนไขคือผู้ให้บริการเครือข่ายทั้งสองฝ่ายต้องเปิด RCS และระบบปฏิบัติการของมือถือต้องใหม่พอ
\end{itemize}

\booksubsection{สถานการณ์การติดตั้ง RCS ในจีน:}

\begin{itemize}
  \item China Mobile: ติดตั้ง RCS ค่อนข้างเร็ว บางมณฑลรองรับบริการ "ข้อความ 5G" แล้ว สามารถส่งข้อความสื่อสมบูรณ์และการ์ดโต้ตอบได้ ต้องมีอุปกรณ์ที่รองรับและเปิดใช้ที่ศูนย์บริการ
  \item China Telecom: ให้บริการผ่าน "China Telecom 5G Message" บางแพ็กเกจรองรับเป็นค่าเริ่มต้น ผู้ใช้เดิมต้องส่งคำสั่งข้อความที่กำหนดเพื่อเปิดใช้
  \item China Unicom: ติดตั้งข้อความ 5G ในเมืองหลัก แต่ความครอบคลุมของอุปกรณ์ไม่สม่ำเสมอ รุ่น Android จำนวนมากยังต้องเปิดสวิตช์ "ข้อความขั้นสูง" หรือ "ข้อความ 5G" ในแอปข้อความด้วยตนเอง
  \item ประสบการณ์จริง: การส่งเสริม RCS ในจีนช้ากว่าที่คาด ส่วนใหญ่ได้รับผลกระทบจากความแข็งแกร่งของระบบนิเวศ WeChat และการปรับให้เข้ากับอุปกรณ์ไม่เป็นมาตรฐานเดียวกัน การสื่อสารในชีวิตประจำวันยังพึ่ง WeChat เป็นหลัก RCS จึงมักทำหน้าที่เป็น "ช่องทางแจ้งเตือนทางการ" (เช่น การแจ้งเตือนจากธนาคาร ข้อความจากรัฐบาล)
\end{itemize}

\begin{tcbwarn}{เส้นแดงเรื่องความปลอดภัยของรหัสยืนยัน}
\begin{itemize}
  \item ห้ามส่งต่อรหัสยืนยัน: ใครก็ตามที่ให้คุณ "อ่านรหัสยืนยัน" หรือ "ส่งต่อข้อความ" อาจเป็นมิจฉาชีพ รวมถึงคนที่อ้างว่าเป็น "ฝ่ายบริการลูกค้า" "พนักงานส่งของ" หรือ "ตำรวจ-อัยการ-ศาล" รหัสยืนยัน = กระเป๋าเงินของคุณ
  \item ระวังข้อความฟิชชิง: ระวังลิงก์ในข้อความเช่น "ETC หมดอายุ" "บัตรประกันสุขภาพถูกอายัด" "พัสดุหาย ขอค่าสินไหม" องค์กรที่ถูกต้องจะไม่ขอเลขบัตรธนาคารหรือรหัสผ่านผ่านลิงก์ในข้อความ
  \item การระบุเบอร์: ข้อความทางการจากธนาคารมักมาจากเลขสั้นขึ้นต้นด้วย 95 หรือ 1069 + ลายเซ็นองค์กร ข้อความ "แจ้งเตือนจากธนาคาร" ที่ส่งจากเบอร์มือถือ 11 หลักทั่วไป แทบทั้งหมดเป็นการหลอกลวง
  \item สำรองการยืนยันสองชั้น: อย่าเปลี่ยนเบอร์มือถือที่ผูกกับบัญชีสำคัญโดยพลการ หากจำเป็นต้องเปลี่ยนเบอร์ ต้องยกเลิกการผูกบัญชีสำคัญทั้งหมดก่อน มิฉะนั้นเจ้าของเครื่องใหม่อาจตั้งรหัสผ่านของคุณใหม่ผ่านรหัสยืนยันทางข้อความ
\end{itemize}

\end{tcbwarn}

\booksection{WeChat/QQ: การจัดการแอปโซเชียลระดับชาติ}

WeChat และ QQ ไปไกลกว่าเครื่องมือสื่อสารแล้ว กลายเป็น “ซูเปอร์แอป” ที่รวมโซเชียล การชำระเงิน และการเข้าถึงข้อมูลไว้ในที่เดียว

\booksubsection{การสำรองและย้ายประวัติแชท:}

\begin{enumerate}
  \item เข้าสู่ระบบ WeChat บนคอมพิวเตอร์ (เวอร์ชัน PC หรือ Mac)
  \item คลิกเมนูมุมล่างซ้าย > [ย้ายและสำรองข้อมูล] > [สำรองและกู้คืน]
  \item เลือก "สำรองประวัติแชทไปยังคอมพิวเตอร์" แล้วเลือกแชทที่ต้องการสำรองบนมือถือ
  \item ย้ายไปมือถือเครื่องใหม่: บนมือถือเครื่องเดิม WeChat [ฉัน] > [การตั้งค่า] > [ทั่วไป] > [ย้ายและสำรองประวัติแชท] > [ย้าย] > [ย้ายไปยังมือถือ/แท็บเล็ต] สแกนคิวอาร์โค้ดก็เสร็จ
\end{enumerate}

\booksubsection{วิธีล้างพื้นที่จัดเก็บแบบเชิงลึก:}

\begin{itemize}
  \item ทางเข้าปกติ: WeChat [ฉัน] > [การตั้งค่า] > [ทั่วไป] > [พื้นที่จัดเก็บ] สามารถล้างแคชและล้างประวัติแชทแบบละเอียดตามเวลาและประเภทไฟล์
  \item ล้างภายในแชท: เข้าแชทใดก็ได้ > "…" มุมขวาบน > [ประวัติแชท] > [จัดการพื้นที่จัดเก็บ] สามารถลบรูป วิดีโอ ไฟล์ในแชทนั้นโดยเฉพาะ โดยเก็บข้อความไว้
  \item ล้างรูป/วิดีโอต้นฉบับ: รูปและวิดีโอความละเอียดสูงต้นฉบับที่รับในแชทกลุ่มกินพื้นที่มากที่สุด สามารถติ๊กเลือก "รูปต้นฉบับ" ใน [จัดการพื้นที่จัดเก็บ] เพื่อลบแบบกลุ่ม
  \item ปิดการดาวน์โหลดอัตโนมัติ: [การตั้งค่า] > [ทั่วไป] > [รูปภาพ วิดีโอ ไฟล์และการโทร] ปิด "ดาวน์โหลดอัตโนมัติ" เพื่อไม่ให้แคชไฟล์กลุ่มจำนวนมากบน WiFi อย่างเงียบ ๆ
  \item แคชมินิโปรแกรม: ใน [ค้นพบ] > [มินิโปรแกรม] กดค้างมินิโปรแกรมที่ไม่ค่อยใช้เพื่อลบ สามารถคืนพื้นที่แคชในเครื่อง
  \item สำรองข้อมูลประวัติแชทแบบบีบอัด: เมื่อสำรองด้วย WeChat บนคอมพิวเตอร์ ติ๊กเลือก "บีบอัด" สามารถบีบอัดประวัติแชทหลาย GB ให้เหลือ 30\%-50\% ของขนาดเดิม
\end{itemize}

\booksubsection{เคล็ดลับการจัดการกลุ่ม:}

\begin{itemize}
  \item ประกาศกลุ่ม: เจ้าของกลุ่ม/ผู้ดูแลโพสต์ใน "…" มุมขวาบนของแชทกลุ่ม > [ประกาศกลุ่ม] จะแจ้งเตือนทุกคนแบบบังคับ เหมาะกับประกาศสำคัญ จะเด้งขึ้นมาอัตโนมัติเมื่อสมาชิกใหม่เข้ากลุ่ม
  \item สิ่งที่ต้องทำในกลุ่ม: เจ้าของกลุ่มสามารถตั้งข้อความปักหมุดเป็น "สิ่งที่ต้องทำของกลุ่ม" สมาชิกทุกคนจะเห็นสัญญาณสีแดงเหนือหัว เมื่อทำเสร็จคลิกแล้วจะหายไปเอง เหมาะกับการรับสมัครกิจกรรมและการเก็บงาน
  \item การต่อแถวในกลุ่ม: พิมพ์ "\#ต่อแถว" ในช่องข้อความเพื่อเริ่มการต่อแถว เหมาะกับการสั่งอาหารและการนับจำนวนสมัคร ผลการต่อแถวจะสรุปรวมที่ด้านบนของแชทกลุ่มอัตโนมัติ ไม่ต้องกังวลว่าแชทจะไหลทับกัน
  \item ไลฟ์ในกลุ่ม: "+" มุมขวาบนของแชทกลุ่ม > [ไลฟ์กลุ่ม] สามารถเริ่มไลฟ์ที่ดูได้เฉพาะในกลุ่มนี้เท่านั้น เหมาะกับการอบรมภายในองค์กร งานเลี้ยงครอบครัว และการทดสอบผลิตภัณฑ์ภายใน
  \item พับเก็บแชทกลุ่ม: ในรายการกลุ่ม ปัดชื่อกลุ่มไปทางซ้าย > [พับเก็บ] สามารถเก็บกลุ่มที่ไม่สำคัญเข้า "แชทกลุ่มที่พับเก็บ" แจ้งเตือนเฉพาะเมื่อมีคน @ คุณเท่านั้น
  \item เฉพาะเจ้าของ/ผู้ดูแลกลุ่มเท่านั้นที่ @ ทุกคนได้: เมื่อกลุ่มมีสมาชิกเกิน 500 คน สมาชิกทั่วไปจะ @ ทุกคนไม่ได้ เพื่อไม่ให้รบกวนกัน
\end{itemize}

\booksubsection{การจัดการสิทธิ์ Moments แบบละเอียด:}

\begin{itemize}
  \item การมองเห็นตามกลุ่ม: เมื่อโพสต์ คลิก "ใครดูได้" เลือก "บางคนเท่านั้น" หรือ "ไม่ให้ใครดู" ตามป้ายกำกับกลุ่มที่ตั้งไว้ล่วงหน้า (เช่น "เพื่อนร่วมงาน" "ครอบครัว" "ลูกค้า") เพื่อส่งถึงคนที่ต้องการอย่างแม่นยำ
  \item ขอบเขตที่อนุญาตให้เพื่อนดู: [ฉัน] > [การตั้งค่า] > [สิทธิ์เพื่อน] > [Moments] > [ขอบเขตที่อนุญาตให้เพื่อนดู Moments] สามารถเลือก "3 วันล่าสุด/1 เดือน/6 เดือน/ทั้งหมด" แนะนำตั้งเป็น "6 เดือน" เพื่อสมดุลระหว่างความเป็นส่วนตัวกับการเข้าสังคม
  \item บล็อก Moments บางโพสต์: กดค้างโพสต์ใน Moments > [ไม่ดู Moments ของเขา/เธอ] หรือ [เลิกติดตาม] อีกฝ่ายจะไม่รู้
  \item จัดการโฆษณาใน Moments: [ฉัน] > [การตั้งค่า] > [เกี่ยวกับ WeChat] > [ข้อมูลส่วนบุคคลและสิทธิ์] > [จัดการโฆษณาส่วนบุคคล] สามารถปิดการแนะนำส่วนบุคคลตามพฤติกรรม
  \item สิบโพสต์สำหรับคนแปลกหน้า: [สิทธิ์เพื่อน] > [Moments] > [อนุญาตให้คนแปลกหน้าดู Moments 10 โพสต์] ปิดได้เพื่อไม่ให้คนแปลกหน้าเลื่อนดูโพสต์ในอดีตของคุณ
\end{itemize}

\booksection{เปรียบเทียบวิดีโอคอล}

เครื่องมือวิดีโอคอลแต่ละตัวมีจุดเด่นต่างกัน การเลือกเครื่องมือที่เหมาะกับสถานการณ์จะทำให้การสื่อสารราบรื่นขึ้น

\begin{itemize}
  \item WeChat วิดีโอ: ครอบคลุมทั่วประเทศ ใช้ได้ทุกวัย เหมาะสำหรับติดต่อเพื่อนฝูงญาติพี่น้องในชีวิตประจำวัน และคุยด่วน รองรับวิดีโอได้สูงสุด 20 คน (ต้องใช้มินิโปรแกรม "วิดีโอกลุ่ม") ข้อเสียคือภาพจะลดคุณภาพชัดเจนเมื่อเครือข่ายไม่ดี
  \item FaceTime: เฉพาะระบบนิเวศ Apple ภาพคุณภาพเยี่ยม เข้ารหัสต้นทางถึงปลายทาง ตั้งแต่ iOS 15 รองรับเชิญผู้ใช้ Android/Web เข้าร่วมผ่าน "ลิงก์" แต่ผู้จัดต้องเป็นอุปกรณ์ Apple เหมาะกับการติดต่อครอบครัวข้ามอุปกรณ์
  \item Zoom: ตัวเลือกแรกสำหรับการประชุมนานาชาติ รองรับการประชุมฟรีสูงสุด 100 คน 40 นาที มีฟังก์ชันครบครันเช่น แชร์หน้าจอ ฉากหลังเสมือน ห้องสนทนาย่อย เหมาะกับการทำงานข้ามประเทศ/ระยะไกล
  \item Tencent Meeting: ตัวเลือกแรกสำหรับงานออฟฟิศในจีน ผสานกับ WeChat อย่างลึกซึ้ง สแกนคิวอาร์โค้ดก็เข้าได้ รองรับการประชุม 300 คน (เวอร์ชันเสียเงิน) บันทึกบนคลาวด์ คำบรรยายสด และสรุปการประชุมมีประสิทธิภาพ เหมาะกับการอบรมภายในองค์กรและการประชุมใหญ่
  \item คำแนะนำตามสถานการณ์: ติดต่อแบบตัวต่อตัวในชีวิตประจำวันเลือก WeChat ระหว่างผู้ใช้ Apple เลือก FaceTime การประชุมธุรกิจอย่างเป็นทางการเลือก Zoom หรือ Tencent Meeting (ในจีนเลือก Tencent ข้ามประเทศเลือก Zoom)
\end{itemize}

\booksubsection{ความต้องการแบนด์วิดท์เครือข่ายและการใช้ข้อมูลอ้างอิง:}

\begin{itemize}
  \item วิดีโอคอล WeChat: ต้องการแบนด์วิดท์ประมาณ 500Kbps-1Mbps คุณภาพมาตรฐานใช้ข้อมูลประมาณ 3-5MB ต่อนาที ความละเอียดสูงประมาณ 8-12MB
  \item FaceTime: ความละเอียดมาตรฐานประมาณ 1Mbps, 720p ประมาณ 2Mbps, 1080p ต้องการ 4Mbps ขึ้นไป ใช้ข้อมูลประมาณ 5-15MB ต่อนาที
  \item Zoom/Tencent Meeting: 720p ประมาณ 1.5-3Mbps, 1080p ต้องการ 3-5Mbps ใช้ข้อมูลประมาณ 300-800MB ต่อชั่วโมง ผู้จัดการประชุมหลายคนจะใช้ข้อมูลมากกว่า
  \item คำแนะนำ: การวิดีโอคอลเป็นเวลานานควรใช้ WiFi บนเครือข่ายมือถือสามารถลดคุณภาพภาพด้วยตนเอง (ปิด "วิดีโอความละเอียดสูง" ในการตั้งค่าแอป) เพื่อประหยัดข้อมูล
\end{itemize}

\booksubsection{เคล็ดลับปรับปรุงคุณภาพภาพวิดีโอคอล:}

\begin{itemize}
  \item แสง: หันหลังให้หน้าต่าง หันหน้าเข้าหาแหล่งกำเนิดแสง (แสงธรรมชาติจะดีที่สุด) หลีกเลี่ยงแสงย้อนที่ทำให้หน้ามืด
  \item มุม: กล้องมือถือควรอยู่สูงกว่าระดับสายตาเล็กน้อย เพื่อไม่ให้เหนียงคางสองชั้น ขณะคอลแนวนอนควรถือสองมือจะนิ่งกว่า
  \item เสียง: ใช้หูฟังเพื่อเลี่ยงเสียงก้องและเสียงหวีด ในที่เงียบอย่าเปิดโหมด "ลำโพง"
  \item เครือข่าย: ใช้ WiFi ก่อน วิดีโอ 4G/5G ควรอยู่ใกล้หน้าต่างเพื่อไม่ให้สัญญาณอ่อน การประชุมหลายคนสามารถปิดวิดีโอของคนอื่น เก็บไว้เฉพาะผู้พูด
  \item ฉากหลังเบลอ: Tencent Meeting และ Zoom รองรับการเบลอโดย AI สามารถพรางฉากหลังที่รกรุงรัง ระวังการแต่งกายสีพื้น อย่าให้สีใกล้กับฉากหลังจนดูเหมือน "หลุดจากกรีนสกรีน"
\end{itemize}

\begin{tcbinfo}{เคล็ดลับเล็ก ๆ}
เมื่อเครือข่ายกระตุก ให้ปิดวิดีโอเก็บเสียงไว้ก่อน สุภาพกว่าการวางสายโดยตรง ก่อนประชุมสำคัญสามารถใช้ฟังก์ชัน “ตรวจเครือข่าย” ของ “Tencent Meeting” ทดสอบแบนด์วิดท์ล่วงหน้า

\end{tcbinfo}

\booksection{อีเมล: การตั้งค่าและการใช้งานบนมือถือ}

\begin{itemize}
  \item เพิ่มบัญชีอีเมล: บน iOS ไปที่ การตั้งค่า > อีเมล > บัญชี > เพิ่มบัญชี บน Android มักมีแอป "อีเมล"
  \item เคล็ดลับการใช้งาน: ตั้งการแจ้งเตือนต่างกันสำหรับบัญชีอีเมลต่าง ๆ ใช้ท่าปัดเพื่อเก็บถาวร ลบ หรือทำเครื่องหมายอีเมลอย่างรวดเร็ว
\end{itemize}

\booksubsection{ความแตกต่างของโปรโตคอล IMAP vs POP3 vs Exchange:}

\begin{itemize}
  \item POP3: หลังดาวน์โหลดอีเมลจากเซิร์ฟเวอร์ จะลบต้นฉบับเป็นค่าเริ่มต้น เก็บไว้ในเครื่องเท่านั้น เหมาะกับการใช้เครื่องเดียวและต้องการอ่านแบบออฟไลน์ ข้อเสียคือเมื่อเปลี่ยนเครื่องจะไม่เห็นอีเมลในอดีต ไม่แนะนำแล้ว
  \item IMAP: อีเมลจะเก็บไว้บนเซิร์ฟเวอร์เสมอ เครื่องเป็นเพียงแคช หลายอุปกรณ์ (มือถือ/คอมพิวเตอร์/แท็บเล็ต) สถานะซิงก์กัน—อ่านที่หนึ่งแล้วถือว่าอ่านทั้งหมด เป็นโปรโตคอลแนะนำหลักในปัจจุบัน
  \item Exchange (รวม ActiveSync): โปรโตคอลของ Microsoft นอกจากอีเมลยังซิงก์ปฏิทิน รายชื่อติดต่อ งาน เป็นตัวเลือกแรกสำหรับอีเมลองค์กร Outlook, NetEase Enterprise Mail และ Tencent Enterprise Mail รองรับทั้งหมด
  \item คำแนะนำในการเลือก: อีเมลส่วนตัวใช้ IMAP อีเมลบริษัทให้ใช้ Exchange ก่อน แล้วค่อย IMAP อย่าใช้ POP3 เว้นแต่ต้องการเก็บถาวรเป็นพิเศษ
\end{itemize}

\booksubsection{วิธีตั้งค่าอีเมลองค์กร:}

\begin{enumerate}
  \item iOS: [การตั้งค่า] > [อีเมล] > [บัญชี] > [เพิ่มบัญชี] > เลือก "Microsoft Exchange" หรือ "อื่น ๆ"
  \item Android: ในแอป "อีเมล" ที่มีในตัว หรือแอป Outlook เลือก "เพิ่มบัญชี" แล้วป้อนที่อยู่อีเมลและรหัสผ่านเต็ม
  \item การกรอกเซิร์ฟเวอร์ (ตัวอย่าง Exchange): ที่อยู่เซิร์ฟเวอร์มักเป็น "mail.โดเมนบริษัท.com" หรือ "outlook.office365.com" โดเมนกรอกชื่อย่อบริษัท
  \item หากแจ้งว่า "ไม่สามารถตรวจสอบเซิร์ฟเวอร์" สามารถป้อนด้วยตนเอง: เซิร์ฟเวอร์ขาเข้า (IMAP) พอร์ต 993 เข้ารหัส SSL เซิร์ฟเวอร์ขาออก (SMTP) พอร์ต 465 หรือ 587 เข้ารหัส TLS
  \item หลังเข้าสู่ระบบ แนะนำเปิด "พุช" (Push) ใน [การตั้งค่าบัญชี] เพื่อรับอีเมลแบบเรียลไทม์ และปิด "ตัวอย่างอีเมล" เพื่อปกป้องความลับทางการค้า
\end{enumerate}

\begin{tcbinfo}{เคล็ดลับเล็ก ๆ}
อีเมลองค์กรมักต้องให้ผู้ดูแลระบบเปิด “รหัสอนุมัติไคลเอนต์” หรือ “รหัสอนุมัติ IMAP/SMTP” ไม่สามารถใช้รหัสผ่านเข้าสู่ระบบเว็บโดยตรงได้ หากเข้าสู่ระบบไม่สำเร็จ ให้ไปสร้างรหัสอนุมัติในการตั้งค่าอีเมลเว็บก่อน แล้วจึงป้อนในมือถือ

\end{tcbinfo}

\begin{tcbpractice}{ลงมือทำ: สำรองประวัติแชท WeChat ไปยังคอมพิวเตอร์สักครั้ง}
\begin{enumerate}
  \item เริ่มสำรองบนคอมพิวเตอร์: เปิด WeChat บนคอมพิวเตอร์ > เมนูมุมล่างซ้าย > [ย้ายและสำรองข้อมูล] > [สำรองและกู้คืน]
  \item เลือกสำรองไปยังคอมพิวเตอร์: คลิก "สำรองประวัติแชทไปยังคอมพิวเตอร์"
  \item เลือกเนื้อหาที่ต้องการสำรองบนมือถือ แล้วรอให้การสำรองเสร็จ
\end{enumerate}

\end{tcbpractice}
